\documentclass[letterpaper,11pt]{article}
\usepackage{times}
\usepackage[margin=1in]{geometry}
\usepackage[singlespacing]{setspace}
\usepackage{amsmath,amssymb}
\usepackage{graphicx,booktabs,array,tabularx}
\usepackage{flafter}
\usepackage{tikz}
\usetikzlibrary{arrows.meta,positioning,fit,automata}
\usepackage[numbers,sort&compress]{natbib}
\usepackage[hidelinks]{hyperref}
\setlength{\emergencystretch}{1.5em}
\newcommand{\researchquestion}[1]{\par\medskip\noindent\textbf{#1}\quad}
\begin{document}
\begin{titlepage}
\begin{center}
\large BOSTON UNIVERSITY\\[.5em]
DEPARTMENT OF ELECTRICAL AND COMPUTER ENGINEERING\\[1.5em]
PhD Prospectus\\[2em]
{\bfseries MULTI-LAYER SECURITY FOR MOBILE USERS:\\
MEASUREMENT AND MITIGATION FROM APPLICATIONS TO CONVERSATION}\\[2em]
By\\[1em]
OLAWALE AMOS AKANJI\\
B.Sc.\ Cyber Security, Air Force Institute of Technology, Nigeria, 2023\\
M.S., Boston University, 2026\\[1.5em]
Advisor: Prof.\ Gianluca Stringhini\\
Co-Advisor: Prof.\ Manuel Egele
\end{center}
\vfill
\begin{center}\bfseries Abstract\end{center}
\noindent
Mobile users can lose control of their data and money through harmful applications,
platform mechanisms that obscure changes in access, and deceptive conversations.
This dissertation investigates these three settings through measurement and the
assessment of practical defenses. Two studies provide its foundation.
\emph{The Cost of Convenience} examines loan applications across five countries
and connects policy requirements to observable application behavior; coordinated
disclosure led to the removal of 93 applications from Google Play.
\emph{Silent Consent, Persistent Risk} studies permission expansion and
cross-application access in Android, combining a longitudinal corpus of
19.3 million APKs with targeted device experiments and a transparency prototype.
The ongoing third study, \emph{SIRENBERT}, investigates whether a sequence of
observable conversational behaviors can support timely warnings about financial
scams on the user's device. An encoder--recurrent-model prototype and a proxy
conversation corpus provide a starting point, while annotation validation,
generalization, warning timing, and deployment evaluation remain research tasks.
The proposed extension examines images and other media presented as evidence of
identity, circumstances, or financial opportunity. I will reconstruct conversations
and embedded images from public screenshot collections, measure how these
artifacts are used, and test whether their interpretation in context improves
warnings over text alone. The central question is how evidence available at each
layer can be translated into useful protection while preserving user privacy and
controlling false alarms. The dissertation will distinguish demonstrated findings
from the hypotheses that these remaining studies will test.
\end{titlepage}

\section{Introduction}
\label{sec:intro}
Mobile devices mediate access to financial services, private communication, and
personal information. Their importance is especially visible where mobile money
provides access to services that conventional banking does not readily
reach~\cite{suri2016mpesa,findex2021}. This dependence creates several opportunities
for abuse. An application may collect information inconsistent with its stated
purpose; a platform may grant access without making the change visible; and a
person using an otherwise legitimate messaging application may be persuaded to
send money under false pretenses. These problems involve different evidence and
intervention points, but share a concern: users may have little opportunity to
recognize a developing risk before a consequential action occurs.

My dissertation examines this concern at the application, platform, and
conversation levels. The first study measures how loan applications comply with
regulatory and platform requirements and how their runtime behavior exposes
sensitive data~\cite{akanji2026cost}. The second investigates Android permission
mechanisms that allow capabilities or cross-application access to expand without
renewed user visibility~\cite{akanji2026silent}. Together, these studies establish
methods for connecting policy or permission abstractions to observable behavior
and for assessing interventions in realistic settings.

The remaining research addresses conversational financial exploitation. Access
controls can constrain an application's use of data, but they do not determine
whether a request made through that application is trustworthy. A request for
money can be legitimate, and behaviors such as affection, professional
self-presentation, and urgency also occur outside scams. The research problem is
therefore to determine whether the evolving context of an exchange provides
useful evidence for a timely warning, without treating ordinary persuasion or
financial discussion as sufficient evidence of fraud.

\textbf{Thesis proposition.} \emph{Measuring how access and trust develop over
time can reveal risks that are obscured by isolated application, permission, or
message observations. These measurements can inform defenses at the point where
the relevant evidence is available, with their value assessed through detection
quality, timeliness, user visibility, and deployment cost.}

This proposition does not require a single detector for all three settings.
The completed studies provide evidence at the application and platform levels.
The ongoing conversation study tests behavioral sequence modeling, and the
proposed extension asks whether images presented within an exchange add useful
evidence. Sections~\ref{sec:loanapps} and~\ref{sec:silentconsent} summarize completed
work; Sections~\ref{sec:sirenbert} and~\ref{sec:proposed} specify what remains to be
established.

\section{Background and Research Gaps}
\label{sec:background}
\subsection{Application behavior and permission visibility}
Digital lending can improve access to credit, but loan applications also create
opportunities for intrusive data collection and coercive debt recovery.
Registration requirements and platform permission policies provide complementary
constraints. The application-level question is whether the behavior of approved
applications is consistent with those requirements, rather than whether approval
alone indicates safety. My loan-app study connects policy interpretation to
static and dynamic measurements~\cite{akanji2026cost}.

Android permissions have been studied in terms of application privileges and
users' understanding of access requests. Felt et al.\ examined permission use
and user attention, while Wijesekera et al.\ investigated the context in which
users make permission decisions~\cite{felt2011demystified,felt2012attention,wijesekera2015remystified}.
Calciati et al.\ measured automatically granted permissions across application
versions; Tuncay et al.\ studied custom-permission weaknesses; and Gamba et al.\
examined preinstalled software~\cite{calciati2020automatically,tuncay2018custom,gamba2020preinstalled}.
My platform study builds on this work by measuring permission evolution at larger
scale and testing selected access paths on Android 16. Its prototype addresses
visibility into changes, rather than claiming to prevent all misuse of an
otherwise granted permission.

\subsection{Conversational scams and evidence beyond text}
Whitty describes the persuasive progression of romance fraud. Oak and Shafiq
use victim interviews to examine relationship-based investment scams, including
trust development, financial exploitation, and subsequent re-engagement.
Acharya and Holz investigate the practices and infrastructure associated with
pig-butchering scams~\cite{whitty2013scammers,oak2025anna,acharya2024explorative}.
These studies motivate examining connected behaviors over time; they do not
imply that every scam follows one fixed sequence or that a particular tactic
uniquely identifies fraud.

Recent detection work provides useful comparison points. Lu et al.\ study mobile
messaging scams using Reddit data, Sun et al.\ introduce a benchmark for scam
progression, and Topcuoglu et al.\ compare generative-model approaches to scam
detection~\cite{lu2026readthis,sun2026prescam,topcuoglu2026genai}.
Their settings and data sources motivate complementary questions about
message-level evidence, transfer between sources, and device constraints.
Our scam-family labels will describe our own data; an external benchmark's
categories are supporting context rather than a required taxonomy.

The remaining gap is the reliability and utility of a representation that can
be updated as a private conversation unfolds. This requires appropriate negative
examples, validated message labels, and evaluation before the financial request.
It also raises a multimodal question: a photograph, document, or balance screenshot
may be presented as proof, while its relevance to a later request depends on what
has already been said. Detecting image reuse or manipulation alone does not
establish that an exchange is exploitative. The proposed work studies the
artifact's role and contextual interpretation together.

\section{Completed Work I: The Cost of Convenience}
\label{sec:loanapps}
\emph{The Cost of Convenience: Identifying, Analyzing, and Mitigating Predatory
Loan Applications on Android}, published at ACM ASIA CCS 2026, investigates
whether loan applications comply with national rules and Google's Financial
Services Policy~\cite{akanji2026cost}.

\textbf{Problem and approach.} An approved lender can still operate an
application that collects prohibited data or exposes it before a borrower has
completed registration. We studied 434 applications from registries and markets
in Indonesia, Kenya, Nigeria, Pakistan, and the Philippines. An LLM-assisted
policy-to-permission mapping translated policy text into testable checks, which
we combined with static code inspection and dynamic observation of network and
data-access behavior.

\textbf{Findings and contribution.} Among approved applications, 141 violated
national regulatory requirements and 147 violated platform policy. In the dynamic
analysis, 37 of 148 functioning applications transmitted sensitive information
at launch after permissions were granted, before signup or a loan application.
Following coordinated disclosure, Google removed 93 flagged applications,
representing over 300 million cumulative installs. These installation counts
indicate potential exposure, not a count of verified victims. The study
contributes a measurement method and evidence for more consistent compliance
assessment across regulatory and platform requirements.

\section{Completed Work II: Silent Consent, Persistent Risk}
\label{sec:silentconsent}
\emph{Silent Consent, Persistent Risk: Android Permission Groups and Custom
Permissions} examines how application capabilities and data access can change
without corresponding user visibility~\cite{akanji2026silent}.

\textbf{Problem and approach.} A permission decision made for one application
version may not convey what a later version can access. Custom permissions can
also protect components in ways that admit unrelated applications. We analyzed
19,324,760 APKs representing 5,965,583 applications in
AndroZoo~\cite{allix2016androzoo}, tracked permission changes, and compared their
association with VirusTotal detections across thresholds~\cite{zhu2020vt}.
Targeted Android 16 experiments tested whether selected mechanisms remained
exploitable. A prototype using public Android APIs exposed otherwise silent
permission changes to the user.

\textbf{Findings and limits.} Of 2,244,575 multi-version applications, 381,026
(17\%) exhibited within-group permission expansion. Malware-flagged applications
had higher odds of expansion at the primary threshold (odds ratio 1.35), an
association rather than evidence that every expansion is malicious. Static
analysis identified 307 active cross-developer access pairs. Dynamic retrieval
of sensitive data was confirmed for the 39 pairs amenable to probing on
Android 16; the dynamic claim does not extend to all 307 pairs. In a 96-day
single-device pilot, the transparency prototype reported 23 expansion events
across 13 applications. This establishes feasibility of providing visibility,
while broader usability and protective effects require further evaluation.

\section{Ongoing Work III: Conversational Scam Detection}
\label{sec:sirenbert}
SIRENBERT investigates whether a mobile device can recognize a developing
conversational scam from the evidence available at each point in an exchange.
The work completed so far comprises a problem definition, a proxy corpus, a
scam-family mapping, a behavioral representation, preliminary measurements,
baseline comparisons, an encoder--recurrent-model prototype, and an Android
integration. The revised annotation and model design described below is the
current specification; its large-scale validation, retraining, generalization
tests, and warning evaluation remain dissertation work.

\subsection{Problem definition and observable decision rule}
A \emph{conversational scam} is a multi-turn exchange in which a suspect, under
false pretenses, induces a target to part with money or other economic value
through a request or instruction. This definition covers romance, investment,
advance-fee, employment, marketplace, impersonation, technical-support, and
related fraud when the financial action is requested through the exchange. The
current evaluated scope requires at least four observed messages and participation
by both parties. It excludes isolated solicitations and coercion-only sextortion,
which do not require the conversational development studied here.

The broad definition describes the phenomenon, but a transcript cannot establish
whether every claimed identity, emergency, or promised return is false in the
world. The operational rule therefore uses observable evidence. A suspect message
is a \emph{financial demand} when it asks or instructs the target to take an action
that would cost money or other economic value. A demand becomes \emph{tactic-linked}
when an earlier or co-occurring tactic in the observed exchange is explicitly
connected to that demand. No completed payment is required. A demand without such
evidence is recorded as \emph{not yet tactic-linked}; this describes what is
observable at that point and does not assert that the request is legitimate.

Let $\mathcal B_t$ be the set of behaviors assigned to suspect message $t$, let
$D$ denote \texttt{Financial\_Demand}, let $\mathcal T$ be the set of tactic
triggers, and let $R(j,t)=1$ when reviewed evidence connects tactic $j$ to demand
$t$. The causal status of a demand is
\begin{equation}
 \ell_t = \mathbb{I}[D\in\mathcal B_t]\,
 \mathbb{I}\!\left[\exists j\leq t:
   \mathcal B_j\cap\mathcal T\neq\varnothing\ \land\ R(j,t)=1\right].
 \label{eq:demand-link}
\end{equation}
Here $\mathbb{I}[\cdot]$ is an indicator: it equals one when the condition in
brackets is true and zero otherwise. Consequently, $\ell_t=1$ only for an
observed demand with at least one reviewed tactic connection available by time
$t$.
The conversation label is binary:
\begin{equation}
 Y(C)=
 \begin{cases}
  \textsc{scam}, & \exists t:\ell_t=1,\\
  \textsc{non-scam}, & \text{otherwise}.
 \end{cases}
 \qquad
 U(C)=\mathbb{I}\!\left[\exists t:
 \mathcal B_t\cap\mathcal T\neq\varnothing\right].
 \label{eq:binary-label}
\end{equation}
Here $U(C)$ records whether a scam-like trajectory is present; it is not a third
label. Thus a conversation containing tactics but no linked demand remains
\textsc{non-scam} with $U(C)=1$. If a first demand is not tactic-linked, later
tactics can still support a warning before a repeated demand; later evidence does
not retroactively change what was observable at the first demand.

The project uses each eligible transcript as observed. A screenshot may begin
mid-conversation, and a complete transcript need not begin with a greeting.
Extraction quality therefore checks whether visible messages were recovered in
order with correct speakers, while capture format is retained as metadata for
stratified analysis. Neither a missing greeting nor source provenance determines
the operational label.

\subsection{Proxy corpus and scam-family coverage}
Complete private scam conversations are scarce and difficult to release. We
assembled a proxy corpus from recorded scam exchanges, public scambait screenshots,
ordinary dialogue, legitimate persuasive requests, and generated conversations.
LoveFraud02 contains researcher-recorded exchanges with purported romance-fraud
perpetrators, while Reddit Scambait contains screenshot-based encounters in which
people knowingly engage suspected scammers~\cite{faber2024lovefraud02,lu2026readthis}.
These sources provide real scam language but are not a representative sample of
unsuspecting victims: scambaiting is performative, screenshot sequences can omit
context, and public reports are selected for posting.

Topical-Chat and XDailyDialog provide ordinary dialogue, and PersuasionForGood
provides legitimate requests for charitable donations
~\cite{gopalakrishnan2019topical,li2017dailydialog,wang2019persuasion}.
The latter supplies a difficult comparison because persuasion and a financial
request are present without fraud. Generated pig-butchering and ScamForge
conversations extend combinations that are rare in the real collections. Their
generation condition is retained as provenance rather than treated as a validated
verdict, and results on real and generated sources will be reported separately.

\begin{table}[tb]
\centering\small
\begin{tabularx}{\linewidth}{@{}Xrr@{}}
\toprule
Source & Conversations & Suspect messages \\
\midrule
LoveFraud02 & 83 & 26,423 \\
Reddit Scambait & 1,166 & 65,685 \\
Topical-Chat & 8,628 & 97,204 \\
XDailyDialog & 3,100 & 19,981 \\
PersuasionForGood & 1,017 & 10,600 \\
Generated pig-butchering scenarios & 5,683 & 237,296 \\
Generated ScamForge scenarios & 3,097 & 152,371 \\
\midrule
Total & 22,774 & 609,560 \\
\bottomrule
\end{tabularx}
\caption{Assembled proxy corpus. Counts describe suspect-side messages available
to the annotation pipeline. Source groups describe how conversations were
obtained; they are not the final transcript-derived labels.}
\label{tab:data}
\end{table}

A separate live and archival Reddit collection has indexed 92,595 distinct
scambait posts and 319,008 Reddit image attachments, most of which are chat
screenshots. The latter is a media-index count that includes unavailable files;
251,624 entries were on disk or recovered at the latest completed audit. A
screenshot may contain multiple visible chat messages, and a post may contain
multiple screenshots, so neither figure represents extracted conversations or
messages. Series reconstruction preserves the
evidence used to join posts that appear to belong to one exchange. These records
are not included in Table~\ref{tab:data}: screenshot extraction, speaker and
message-order checks, deduplication, and a held-out quality audit must precede
their admission to the conversation corpus. The same image-preserving collection
also supplies the primary candidate pool for the proposed multimodal study in
Section~\ref{sec:proposed}. It retains the source images needed to recover both
message text and embedded credibility artifacts such as identity documents,
receipts, photographs, and balance or investment screenshots.

To describe coverage without conflating a scam name with its mechanism, the
current mapping follows seven expected-benefit or expected-consequence families
from Beals et al.~\cite{beals2015fraud}. Each conversation receives one primary
family; relationship use, impersonated identity, payment mechanism, and a common
name are orthogonal modifiers. For example, pig butchering that promises returns
belongs to Consumer Investment, with \texttt{relationship\_used=true} and
\texttt{named\_scam=pig\_butchering}. ``Advance fee'' records how value is taken
and can occur in several families, so it is not a universal primary family.

\begin{table}[tb]
\centering\small
\begin{tabularx}{\linewidth}{@{}p{.28\linewidth}XX@{}}
\toprule
Primary family & Organizing expectation & Examples or modifiers \\
\midrule
Consumer Investment & Return on an investment & crypto, forex, investment-return pig butchering \\
Consumer Products/Services & Product, service, credit, or sale & marketplace, fake buyer, tech support, recovery, fake loan, paid service \\
Employment & Income from work & fake job, reshipping, mystery shopper, multilevel marketing \\
Prize/Grant & Prize, grant, inheritance, or windfall & lottery, sweepstakes, inheritance, government grant, Nigerian letter \\
Phantom Debt Collection & Avoiding a claimed debt consequence & tax, toll, court, or medical debt \\
Charity & Supporting a claimed cause or hardship & disaster, religious, personal appeal \\
Relationship/Trust & Maintaining or helping a relationship & romance, allowance, grandparent, friend/relative, boss request \\
\bottomrule
\end{tabularx}
\caption{Conversation-level family mapping. These labels describe corpus
coverage and support family-held-out evaluation; the completed-conversation
family label is not supplied to the causal detector.}
\label{tab:families}
\end{table}

An initial model-assisted pass over 520 in-scope real-source candidates proposed
218 Relationship/Trust, 78 Products/Services, 53 Prize/Grant, 31 Investment,
13 Employment, eight Charity, and two Phantom Debt cases; it also returned 115
\texttt{NONE} and two unparsed outputs. That distribution is provisional. An
automated audit flagged 67 of the 520 outputs. Independent review of the 115
\texttt{NONE} cases agreed on the apparent family in 97 cases and assigned 108 to
one of the seven families, showing that the unreviewed model output substantially
undercounted several families. The final distribution will be reported only after
all flagged and probability-sampled cases receive a second reading.

\subsection{Lifecycle groups, triggers, and forms}
The representation separates a conversation family from what each message does.
It uses four non-exclusive evidence groups:
\begin{equation}
 \text{Contact}\ \longrightarrow\ \text{Relationship}\
 \ \longrightarrow\ \text{Tactics}\ \rightleftarrows\
 \ \text{Financial Demand}.
 \label{eq:lifecycle}
\end{equation}
The arrows summarize common progressions rather than impose a required order.
Groups may be skipped, repeated, or revisited after a demand. In particular,
Relationship is not a benign state, Tactics is one group rather than three states,
and Financial Demand is an observable act rather than an automatic verdict.

The four groups are separated by the function a behavior serves in relation to
financial extraction, as summarized in Table~\ref{tab:lifecycle-functions}.

\begin{table}[tb]
\centering\small
\begin{tabularx}{\linewidth}{@{}p{.20\linewidth}X@{}}
\toprule
Group & Function in the observed conversation \\
\midrule
Contact & Opens the interaction. \\
Relationship & Builds or maintains the interaction itself. \\
Tactics & Creates a reason, condition, or pressure for the target to give up economic value. \\
Financial Demand & Asks the target to give up that value. \\
\bottomrule
\end{tabularx}
\caption{Functional definition of the four lifecycle groups. The groups describe
what a behavior does, rather than require it to occur at a fixed point in time.}
\label{tab:lifecycle-functions}
\end{table}

The central boundary is between Relationship and Tactics. We apply a functional
test: could the observed behavior serve as the reason, condition, or lever for a
payment request? Routine conversation, personal disclosure, and affection build
the interaction but do not, by themselves, supply that basis. A claimed identity
or status, promised return, crisis, verification avoidance, request for secrecy,
deadline, or emotional condition can do so. This division parallels prior
lifecycle accounts that distinguish grooming from the sting and bonding from the
bait or squeeze~\cite{whitty2013scammers,oak2025anna}.

Romantic affection therefore remains a Relationship behavior even when it occurs
inside a romance scam. For example, ``I love you'' does not itself make a later
request for money tactic-linked. A fabricated hospital crisis, or making affection
conditional on payment, supplies Tactics evidence that may support such a link.
Operationally, Relationship evidence alone cannot satisfy the demand-link relation
in Equation~\ref{eq:demand-link}. Tactics evidence is eligible to do so only when
the annotation records an explicit connection to the demand. A demand for which no
such connection is yet visible is \texttt{not\_yet\_tactic\_linked}; this describes
the observed evidence and does not establish that the request is legitimate.

\begin{figure}[tb]
\centering\resizebox{\linewidth}{!}{\begin{tikzpicture}[font=\scriptsize,>=Latex,
 state/.style={draw=blue!55!black,rounded corners,align=center,minimum height=.8cm,fill=blue!4,text width=2.05cm},
 arr/.style={->,draw=blue!55!black,thick}]
\node[state] (n) at (0,0) {Start\\$s=0$};
\node[state] (i) at (3,0) {Initial\\contact\\$s=1$};
\node[state] (r) at (6,0) {Relationship\\building\\$s=2$};
\node[state] (p) at (9,0) {Pre-exploitation\\behavior\\$s=3$};
\node[state] (e) at (12,0) {Demand after\\setup\\$s=4$};
\draw[arr] (n)--node[above=4mm]{contact}(i);
\draw[arr] (i)--node[above=4mm]{rapport}(r);
\draw[arr] (r)--node[above=4mm]{setup}(p);
\draw[arr] (p)--node[above=4mm]{demand}(e);
\draw[arr] (n.north) to[bend left=25] node[above]{setup may appear immediately} (p.north);
\draw[arr] (i.south) to[bend right=23] node[below]{setup may precede rapport} (p.south);
\draw[arr] (p) edge[loop above,looseness=5,min distance=11mm] node[above]{further setup / ordinary talk} (p);
\draw[arr] (e) edge[loop right,looseness=5,min distance=8mm] node[right]{later messages} (e);
\end{tikzpicture}
}
\caption{Current lifecycle representation. The detector consumes the ordered
message sequence; the four groups organize evidence for interpretation and do
not force a monotone state path.}
\label{fig:lifecycle}
\end{figure}

The candidate annotation vocabulary contains thirteen observable triggers
(Table~\ref{tab:triggers}). Every suspect message receives at least one trigger,
and independently supported triggers can co-occur. \emph{Routine Engagement}
labels a greeting, acknowledgment, check-in, joke, or everyday exchange that
keeps an established conversation active when no more specific behavior is
expressed. It is therefore a positive account of ordinary maintenance rather
than an ``other'' category. The first visible message is assigned
\emph{Contact Initiation} only when its content or visible context actually
establishes the interaction.

\begin{table}[tb]
\centering\scriptsize
\begin{tabularx}{\linewidth}{@{}p{.13\linewidth}p{.23\linewidth}Xp{.25\linewidth}@{}}
\toprule
Group & Trigger & Observable function & Closed forms, when used \\
\midrule
Contact & Contact\_Initiation & Attempts to establish the interaction & ordinary greeting; wrong number; mistaken identity; direct approach \\
\midrule
Relationship & Routine\_Engagement & Keeps an established interaction active & --- \\
& Personal\_Disclosure & Shares or elicits personal, vulnerable, or intimate information & --- \\
& Affectionate\_Bonding & Builds or intensifies affection or romantic commitment & warmth/endearment; romantic interest; commitment \\
\midrule
Tactics & Credential\_Positioning & Invokes identity, standing, resources, expertise, access, or connections for credibility & identity/relationship; authority/status; wealth/resources; profession/expertise; access/connection \\
& Gain\_Offer & Presents an economic opportunity or benefit & investment return; employment income; prize/grant/inheritance; loan/credit/debt relief; gift/allowance; product/service; donation/support \\
& Crisis\_Disclosure & Presents an emergency or acute hardship & --- \\
& Verification\_Avoidance & Evades an independent check or substitutes controlled proof & --- \\
& Platform\_Migration & Requests movement to another communication channel & --- \\
& Secrecy\_Request & Requests concealment or deletion of evidence & --- \\
& Urgency\_Pressure & Uses a deadline, scarcity, or demand for immediate action & --- \\
& Emotional\_Leverage & Makes affection, guilt, obligation, approval, fear, or relational loss conditional on compliance & --- \\
\midrule
Demand & Financial\_Demand & Requests or instructs an action costing economic value & payment method recorded separately \\
\bottomrule
\end{tabularx}
\caption{Candidate multilabel trigger vocabulary and closed forms. Forms preserve
distinctions within a shared conversational function without turning each
surface realization into another top-level trigger.}
\label{tab:triggers}
\end{table}

This granularity follows a functional rule: a top-level trigger must change what
the sequence model should retain about the interaction. Closely related content
with the same function remains a form. Thus a claim of authority and a claim of
wealth are distinct Credential Positioning forms because both attempt to raise
the standing on which later claims are believed; the planned ablation will test
whether retaining those forms improves detection. The inventory claims coverage
of observable behaviors relevant to the in-scope scam definition, not every act
in human conversation. Its coverage will be tested by residual review across all
seven families, explicit searches for missed acts, co-occurrence analysis, and
inter-annotator agreement. A new trigger is added only when reviewed examples
show an independently observable function that the existing triggers cannot
represent and that Stage~2 can use.

For interpretation, define a four-dimensional accumulated-presence summary
\begin{equation}
 \mathbf s_t=[c_t,r_t,u_t,d_t],\qquad
 \mathbf s_t=\mathbf s_{t-1}\lor G(\mathcal B_t),
 \label{eq:evidence-summary}
\end{equation}
where $G$ maps the current behavior set to Contact, Relationship, Tactics, and
Financial Demand, and $\lor$ is element-wise logical OR. This summary permits
skipped groups but loses repetitions after a bit first turns on. The learned
model therefore receives the full ordered per-message sequence; Equation
~\ref{eq:evidence-summary} is used only for explanation or as an ablation.

For a financial demand, annotation records only information that has a defined
role: one or more payment methods and the evidence links used to validate
Equation~\ref{eq:demand-link}. The closed payment list is cash, bank or wire,
card, gift card, cryptocurrency, peer-to-peer application, cheque, online
platform, mobile money or airtime, goods or property, and unspecified. Requested
amount and recipient are not mandatory structured fields: an exploratory scan
found amount mentions at similar rates in scam and benign messages, while those
values can still be extracted later for descriptive analysis.

\subsection{Preliminary measurements and baselines}
The earlier single-label annotations provide motivation, not final ground truth.
Under their automated rule, 520 real-source conversations contained a qualifying
demand. They had a median of two distinct mapped indicators before the first
demand, and 52.3\% had at least two. Among 1,685 Topical-Chat, XDailyDialog, and
PersuasionForGood conversations containing a demand but not receiving that scam
label, 90.1\% contained no mapped indicator before the demand. Indicator frequency
also varied with length: a money-making offer appeared in 33.6\% of the shortest
quartile but 13.2\% of the longest, while affection and emotional leverage became
more common in longer conversations. This is consistent with shorter investment
approaches and longer relationship approaches, and cautions against treating one
tactic as necessary.

A conversation-level audit selected one demand from each of 751 rule-positive and
1,337 rule-negative conversations in the real-source subset. The first selected
demand occurred within the first four suspect messages in 53 of 751 positives
(7.1\%) and 735 of 1,337 negatives (55.0\%). Gift-card, cryptocurrency, or
peer-to-peer payment language appeared in 98 positives and eight negatives; among
the early demands it appeared in nine positives and five negatives. These values
show why payment method is a candidate feature, but they do not estimate field
precision. The benign sources contain little realistic person-to-person payment
conversation, so payment-method gains must be tested on a dedicated benign payment
stress set and under source-held-out evaluation.

We also compared demand-only and completed-conversation baselines over all 22,774
conversations. For this preliminary comparison only, source provenance supplies
the positive and negative groups because the revised transcript labels are not
yet validated. The five encoders were trained previously to recognize heuristic
financial-demand labels; the zero-shot Gemma 4 E4B judge received the completed
conversation. Table~\ref{tab:baselines-22774} therefore measures separation of
the assembled source collections rather than performance against the final
operational rule.

\begin{table}[tb]
\centering\scriptsize
\setlength{\tabcolsep}{3.5pt}
\begin{tabular}{@{}lcrr@{}}
\toprule
System & Scam $F_1$ [95\% CI] & Scam-source missed & Benign-source flagged \\
\midrule
Keyword financial-demand rule & .718 [.709, .727] & 797 & 2,529 \\
DistilBERT demand classifier & .710 [.701, .719] & 789 & 2,676 \\
BERT-base demand classifier & .712 [.703, .721] & 791 & 2,634 \\
FinBERT demand classifier & .713 [.703, .722] & 795 & 2,620 \\
DeBERTa-v3 demand classifier & .705 [.695, .714] & 842 & 2,667 \\
MobileBERT demand classifier & .716 [.706, .725] & 808 & 2,546 \\
Zero-shot completed-conversation judge & .896 [.890, .903] & 296 & 797 \\
\bottomrule
\end{tabular}
\caption{Preliminary source-proxy baselines. Confidence intervals use 5,000
multinomial bootstrap draws. The completed-conversation judge can use evidence
after a demand and therefore is not an early-warning baseline.}
\label{tab:baselines-22774}
\end{table}

The demand detectors cluster between $F_1=.705$ and $.718$, supporting the
motivation that recognizing a request for money does not determine whether it is
fraudulent. The completed-conversation judge reaches $F_1=.896$ but still misses
296 scam-source conversations and flags 797 benign-source conversations; it also
has access to future messages unavailable to a causal detector.

The implemented encoder--GRU prototype supplies a second starting point. Nested
five-fold cross-validation selected ModernBERT with ten preceding messages. On
the original 1,297-conversation test partition and its original labels, it
obtained scam $F_1=.863\pm.005$ across seeds 7, 42, and 123, recall from .875 to
.917, and eight to ten false positives among 1,249 original negatives. Per-seed
bootstrap intervals spanned [.775, .933]. Re-scoring the unchanged models against
the provisional revised labels reduced mean $F_1$ to $.701\pm.007$ and mean recall
to .610 because the positive count changed from 48 to 70. This is evidence of
label mismatch, not a retrained result; the revised annotation must be validated
before the model is retrained and a new headline value is reported.

\subsection{Revised two-stage model}
\textbf{Stage 1: behavioral interpretation.} Let $x_{1:t}$ denote the current
message and all context from both speakers available through that message.
DistilBERT, ModernBERT,
and an adapted on-device language model will be compared as contextual encoders
~\cite{sanh2019distilbert,warner2024modernbert,hu2022lora}. Each encoder produces
a shared representation
\begin{equation}
 \mathbf h_t^{\mathrm{enc}}=\operatorname{Encoder}_{\theta}(x_{1:t}).
 \label{eq:shared-encoder}
\end{equation}
Thirteen independent sigmoid outputs preserve co-occurring behaviors:
\begin{equation}
 \mathbf p_t^{\mathrm{trigger}}=
 \sigma(\mathbf W_{\mathrm{trig}}\mathbf h_t^{\mathrm{enc}}+
 \mathbf b_{\mathrm{trig}})\in[0,1]^{13}.
 \label{eq:trigger-head}
\end{equation}
These values need not sum to one. A softmax would force one label and reproduce
the loss of information observed when an urgent gift-card demand could be labeled
only as Financial Demand.

Four conditional form heads represent four Contact forms, three Affectionate
Bonding forms, five Credential Positioning forms, and seven Gain Offer forms.
For parent trigger $r$,
\begin{equation}
 \mathbf q_t^{\mathrm{form}(r)}=
 \operatorname{softmax}(\mathbf W_r\mathbf h_t^{\mathrm{enc}}+\mathbf b_r),
 \qquad
 \mathbf p_t^{\mathrm{form}(r)}=p_t(r)\mathbf q_t^{\mathrm{form}(r)}.
 \label{eq:form-head}
\end{equation}
The four gated blocks concatenate to
$\mathbf p_t^{\mathrm{form}}\in[0,1]^{19}$; each form loss is evaluated only when
its parent trigger is present. Payment methods are multilabel because one demand
can name more than one mechanism:
\begin{equation}
 \mathbf q_t^{\mathrm{pay}}=
 \sigma(\mathbf W_{\mathrm{pay}}\mathbf h_t^{\mathrm{enc}}+\mathbf b_{\mathrm{pay}}),
 \qquad
 \mathbf p_t^{\mathrm{pay}}=p_t(D)\mathbf q_t^{\mathrm{pay}}
 \in[0,1]^{11}.
 \label{eq:payment-head}
\end{equation}
The value \texttt{unspecified} is positive only when a demand names none of the
ten explicit methods. Payment loss is evaluated only for demand messages.

An optional learned projection retains contextual wording outside the structured
inventory:
\begin{equation}
 \mathbf e_t=\mathbf W_{\mathrm{proj}}\mathbf h_t^{\mathrm{enc}}+
 \mathbf b_{\mathrm{proj}}\in\mathbb R^d,
 \qquad
 \mathbf z_t=[\mathbf p_t^{\mathrm{trigger}};
 \mathbf p_t^{\mathrm{form}};\mathbf p_t^{\mathrm{pay}};\mathbf e_t]
 \in\mathbb R^{43+d}.
 \label{eq:stage1-vector}
\end{equation}
Stage~2 receives this one concatenated numeric vector, not raw text and not gold
annotations. With $d=64$, for example, $\mathbf z_t$ has 107 dimensions.

\textbf{Stage 2: causal aggregation.} A unidirectional GRU processes the ordered
Stage~1 vectors~\cite{cho2014gru}:
\begin{equation}
 \mathbf g_t=\operatorname{GRU}_{\phi}(\mathbf z_t,\mathbf g_{t-1}),
 \qquad
 \hat y_t=\sigma(\mathbf w^{\top}\mathbf g_t+b).
 \label{eq:stage2}
\end{equation}
The recurrent memory can represent order, repetition, co-occurrence, demand
position, and tactics that appear between repeated demands. The final family,
demand links, demand-link status, and conversation label are supervision or
evaluation outputs and are never detector inputs.

\begin{figure}[tb]
\centering\resizebox{\linewidth}{!}{\begin{tikzpicture}[font=\small, >=Latex,
 block/.style={draw=blue!55!black, rounded corners=2pt, fill=blue!4, align=center, minimum height=1.18cm},
 flow/.style={->,thick,draw=blue!55!black}]
\node[block,text width=2.6cm] (input) at (0,0) {Message $m_t$\\+ recent context};
\node[block,text width=3.0cm] (s1) at (3.75,0) {\textbf{Stage 1}\\Behavior classifier\\$\mathbf z_t\in[0,1]^{14}$};
\node[block,text width=3.0cm] (s2) at (7.85,0) {\textbf{Stage 2}\\Sequence model\\memory $\mathbf h_t$};
\node[block,text width=2.2cm] (warning) at (11.7,0) {Risk $r_t$\\Warning policy};
\draw[flow] (input)--(s1);
\draw[flow] (s1)--node[above,font=\scriptsize]{$\mathbf z_t$}(s2);
\draw[flow] (s2)--(warning);
\draw[flow] (s2.south east) -- ++(0,-.42) -| (s2.south west);
\node[font=\scriptsize,align=center] (memory) at (7.85,-1.35) {Memory carried to the next message};
\node[draw=blue!55!black,dashed,rounded corners,fit=(input)(s1)(s2)(warning)(memory),inner xsep=7pt,inner ysep=13pt] (device) {};
\node[anchor=south west,font=\scriptsize\bfseries,text=blue!55!black] at (device.north west) {ON-DEVICE DETECTION: TEXT PROTOTYPE AND PLANNED MODEL COMPARISONS};
\end{tikzpicture}
}
\caption{Revised two-stage architecture. Stage~1 predicts behaviors and selected
attributes from the available prefix. Stage~2 aggregates only causal numeric
outputs and produces an updated risk score after each message.}
\label{fig:architecture}
\end{figure}

The design will be tested by four nested ablations: triggers alone; triggers plus
forms; triggers, forms, and payment method; and the full vector with the compact
text projection. The payment block will be retained only if gains survive testing
on realistic benign payment conversations and source-held-out splits. The text
projection can also encode source style, so it will be retained only if its gains
survive both source-held-out and family-held-out evaluation. A random-split gain
alone is insufficient.

\subsection{Validation and remaining work}
The next annotation pilot will compare the earlier single-label output with the
candidate multilabel representation on a coverage-based sample containing
positive and boundary cases for every trigger, form, payment method, and important
co-occurrence, across transcript and screenshot sources. At least two independent
annotators will label the same conversations. Validation will report prevalence,
positive agreement, and a chance-corrected agreement statistic for each trigger;
kappa alone is insufficient for rare labels. Reviewers will also search for
relevant behaviors outside the inventory. Only after semantic disagreements and
evidence-link errors are resolved will the full corpus be annotated.

\researchquestion{RQ-T: Can a causal behavioral sequence support timely and
transferable scam warnings?}
Training, validation, and test partitions will keep related posts, fragments, and
generated variants together~\cite{cawley2010overfitting,varma2006bias}. Evaluation
will report message-level trigger and attribute accuracy, conversation-level
precision and recall, macro-$F_1$, calibration, and false alerts per conversation.
Entire-source and leave-one-family-out partitions will test whether the model
learns the intended behavior sequence rather than source style or one scam family.

Warning time will be measured causally at four checkpoints: before the first
observed demand, at that demand, after it but before a repeated demand, and by the
end of the observed exchange. If $d_1$ is the first tactic-linked demand and
$\tau$ is selected on validation data, define
\begin{equation}
 \hat t_\tau=\min\{t:\hat y_t\geq\tau\},\qquad
 L_\tau=d_1-\hat t_\tau.
 \label{eq:lead}
\end{equation}
Positive $L_\tau$ is a warning before the demand. Missed warnings will be reported
separately rather than omitted from average lead time. At a detected financial
demand below the scam threshold, the product can present a neutral pause-and-verify
notice without claiming that the request is fraudulent.

Adversarial evaluation will insert plausible ordinary messages before or between
tactics, because padding can push evidence outside a bounded context or dilute a
learned sequence. Paraphrased tactics, delayed and repeated demands, changed
payment methods, and removed platform-specific terms will test related failure
modes. On-device replay will measure agreement with the research model,
quantization effects, latency, peak memory, and energy on named hardware. The
Android integration establishes an implementation path; the annotation study,
retraining, holdout experiments, early-warning evaluation, and device measurements
are the work required to establish the proposed system.


\section{Proposed Work: Credibility Artifacts in Conversational Scams}
\label{sec:proposed}
\subsection{Problem and motivation}
A conversational claim may be supported by a photograph, identity document,
receipt, or investment screenshot. I use \emph{credibility artifact} to mean
media presented within an exchange as evidence for an identity, circumstance,
capability, transaction, or opportunity. The term describes its use; it does not
presume that the artifact is forged or malicious. Ordinary photographs and
documents may serve the same function in legitimate exchanges.

The unresolved problem is whether these artifacts provide reliable additional
evidence for scam detection, and whether that evidence arrives early enough to
improve warnings. A text-only model may miss what a screenshot claims; an image
classifier may miss why that claim matters to the request. Conversely, an
apparently inconsistent or reused image may be benign, and a scammer can use
genuine media. The proposed study will measure these cases before deciding how
to incorporate artifact evidence into the detector. Its primary empirical scope
is images and screenshots. Audio will be a conditional extension if usable
original recordings become available; an audio-message icon in a screenshot is
not an audio recording.

\begin{figure}[htbp]
\centering\resizebox{\linewidth}{!}{\begin{tikzpicture}[font=\small,>=Latex,
 b/.style={draw=teal!65!black,fill=teal!4,rounded corners=2pt,align=center,minimum height=1.05cm,text width=3.35cm},
 a/.style={->,thick,draw=teal!65!black}]
\node[b] (raw) at (0,0) {Public screenshot posts\\Images and provenance};
\node[b] (rec) at (5,0) {Reconstruct exchange\\Text, speaker, order, crops};
\node[b] (valid) at (10,0) {Quality review\\Accept, correct, or withhold};
\node[b] (measurement) at (0,-2) {\textbf{RQ1: Measurement}\\Artifact type, timing, reuse};
\node[b] (claims) at (5,-2) {\textbf{RQ2: Interpretation}\\Claims and evidence\\in context};
\node[b] (fusion) at (10,-2) {\textbf{RQ3: Detection}\\Incremental value on device};
\draw[a] (raw)--(rec); \draw[a] (rec)--(valid);
\draw[a] (valid.south)--(fusion.north);
\draw[a] (valid.south) -- ++(0,-.33) -| (measurement.north);
\draw[a] (measurement)--(claims); \draw[a] (claims)--(fusion);
\end{tikzpicture}
}
\caption{Proposed artifact study. Collection and extraction supply candidate
evidence; quality review precedes measurement and detector evaluation. Artifact
presence and source membership do not supply automatic scam labels.}
\label{fig:artifacts}
\end{figure}

\researchquestion{RQ1: What artifacts appear, and how are they used over a conversation?}
\textbf{Why this matters.} A detector design needs evidence about which media
are available, what claims they support, and when they appear. Without this
measurement, a multimodal extension could focus on visually distinctive but
rare or late-arriving content.

\textbf{Approach.} I will use the existing screenshot collection and the completed
live and historical Reddit post index described in Section~\ref{sec:sirenbert}.
Its full-resolution images support both conversation-text reconstruction and the
recovery of media embedded within those conversations. The collection retains
post identifiers, gallery order, available media, and retrieval status; media
backfill and conversation extraction remain in progress.
I will reconstruct messages and embedded artifact regions while retaining
links to their source pixels. Text recognition, spatial layout, and a vision
language model will be compared for transcription and speaker attribution.
Different extraction systems' agreement will prioritize review, but a human
sample of agreeing outputs is also necessary to detect shared mistakes.

Posts will be joined into a conversation only when continuity evidence supports
the connection; matching an author or title is not sufficient. Screenshot overlap
will be removed without deleting genuinely repeated messages. The review will
distinguish the screenshot owner's messages from those of the suspected scammer,
retain uncertain roles, and record missing sections. I will audit extraction
quality by app interface, image length, compression, and source before measuring
artifact prevalence, first appearance, and within/across-conversation reuse.
Exact and perceptual image matches will identify reuse candidates, followed by
review. A match demonstrates shared content, not identity theft or common criminal
control. Synthetic-image detectors may provide exploratory scores, but those
scores will not be reported as the prevalence of AI-generated scam imagery
without suitable ground truth and calibration.

\textbf{Evaluation and expected contribution.} The study will report message
coverage, transcription error, speaker accuracy, artifact crop/link accuracy,
and extraction yield before giving corpus-level measurements. A held-out audit
will quantify selection and reconstruction limitations. The contribution will
be a provenance-preserving measurement corpus and evidence about artifact use
in observed exchanges, not a claim to represent all conversational fraud.

\researchquestion{RQ2: Does an artifact support or conflict with the claims made in context?}
\textbf{Why this matters.} Recognizing a face or extracting a balance does not
establish what the sender asserts. A uniformed photograph is inconsistent with
an occupation claim only if the sender presents the depicted person as themselves
and the context makes those claims incompatible. The study must distinguish
contradiction from irrelevance, ambiguity, and missing evidence.

\textbf{Approach and evaluation.} I will annotate the artifact's visible content,
the claim explicitly associated with it, and the earlier messages needed to
interpret that claim. Candidate relations will include support, contradiction,
and insufficient information. Text-only, artifact-only, and joint models will
be compared on reviewed examples, including genuine documents, benign image
reuse, edited screenshots, and scams using apparently consistent media.
Controlled mismatches will test specific capabilities but will be reported
separately from naturally occurring cases. Evidence localization and
claim-relation accuracy will accompany classification metrics. The expected
contribution is an evaluated contextual-evidence representation, including its
failure and abstention cases, rather than automatic verification of identity.

\researchquestion{RQ3: Does artifact evidence improve warnings under mobile constraints?}
\textbf{Hypothesis.} When text alone leaves a claim ambiguous, an artifact's
content and its relation to prior messages may improve detection at an acceptable
false-alarm rate. Additional model complexity is justified only if the measured
benefit survives source separation and device constraints.

I will compare simple score combination with learned fusion. Let $\mathbf e_t$
represent the artifact evidence available at message $t$, and let $u_t$ indicate
whether an artifact is available. A candidate extension of Stage 2 is
\begin{equation}
\mathbf h_t^{\mathrm{multi}}=
 g_\psi\!\left(\mathbf h_{t-1}^{\mathrm{multi}},
 [\mathbf z_t;\mathbf e_t;u_t]\right).
\label{eq:fusion}
\end{equation}
When no artifact is present, $u_t=0$ and a masked representation supplies no
invented evidence. This formulation does not require turning every image type
into a new trigger. The experiment will determine whether the representation
and fusion mechanism help beyond the existing behavioral sequence.

Evaluation will compare text alone, artifacts alone, and the combined system on
the same eligible conversations, with component ablations. Related conversations
and reused-image clusters will be separated across evaluation partitions to
limit memorization. I will report recall at matched false-alarm rates, warning
lead time, calibration, and robustness to cropping, compression, missing media,
and benign media inserted to distract the detector. Device evaluation will include
latency, energy, memory, and the cost of retaining context. The default deployment
design will process private media locally. Any remote lookup would require its
own privacy and utility analysis; privacy does not follow from using an image hash.

\subsection{Feasibility, risks, and completion criteria}
A live and archival Reddit collection is complete, while media backfill and the
screenshot-extraction prototype remain underway. Initial extraction runs
demonstrate that structured text and artifact candidates
can be recovered, while tall screenshots, speaker confusion, and tile overlap
remain concrete challenges. These are feasibility observations, not a completed
accuracy evaluation. Screenshot crops may be too small or degraded for reliable
artifact analysis; the study will report that yield and can still characterize
observable use and timing without claiming to determine authenticity.

The first milestone is a reviewed reconstruction sample and a frozen extraction
contract. The second is a corpus measurement with documented selection and error
rates. The third is a held-out comparison establishing when artifact evidence
helps, harms, or adds no measurable value. A negative detection result remains
informative if the evaluation identifies why the evidence is unavailable,
unreliable, or redundant. Before release, personal identifiers will be reviewed
and redacted where appropriate; data availability will reflect source permissions
and human-subjects requirements rather than assume all collected media can be
redistributed.

\section{Research Plan and Expected Contributions}
\label{sec:timeline}
The remaining plan prioritizes reliable evidence before new model complexity.
The first twelve months focus on the conversation-label pilot, reconstruction
quality, text-model evaluation, and artifact measurement. The timetable below
is provisional and depends on review capacity and the yield of usable media.

\begin{table}[htbp]
\centering\small
\begin{tabularx}{\linewidth}{@{}p{.19\linewidth}X@{}}
\toprule
Period & Deliverable and decision point \\
\midrule
Fall 2026 & Prospectus defense; freeze and review the annotation guide; validate
reconstruction on held-out screenshots; establish an eligible corpus and leakage-aware splits. \\
Spring 2027 & Retrain/evaluate text models; measure warning timing, source/family
transfer, and padding robustness; complete initial artifact measurements. \\
Summer 2027 & Complete the contextual-claim reference set and compare interpretation
methods; decide which artifact evidence merits detector integration. \\
Fall 2027 & Evaluate fusion and device costs; investigate warning presentation and
cross-app continuity within approved study scope. Prepare manuscripts. \\
Spring 2028 & Resolve remaining evaluations, document limitations and artifacts,
and complete dissertation writing and defense. \\
\bottomrule
\end{tabularx}
\caption{Provisional milestones. Model expansion follows evidence and quality
checks rather than a predetermined accuracy or publication outcome.}
\end{table}

The completed studies contribute methods and findings about application compliance
and permission visibility. The ongoing study is expected to contribute a validated
behavioral representation and an evaluation of sequential, on-device scam
warnings. The proposed artifact study will add a contextual measurement resource
and evidence about the value and limits of multimodal detection. Together, these
contributions will connect mobile-security measurements to interventions whose
benefits and limitations are explicitly tested.

{\footnotesize\bibliographystyle{unsrtnat}\bibliography{prospectus_revised_draft}}
% Attach the candidate's current CV separately if required by the department.
\end{document}
